\section{Related work}
\label{sec:related}

\paragraph{Jev-style decision models.}
Most of the community decision models compared in LLM-as-Jev \citep{li2026llmasjev} use a decoder backbone and read the logits of an option letter or label, with the backbone frozen or tuned with LoRA.
LLM-as-Jev studies this readout and when fine-tuning helps.
AnyJev reads a typed decision from one prefill and corrects the label prior and the effect of option order without training \citep{zhang2026anyjev}.
The Contrastive Language Model \citep{kwok2026clm} embeds the state (the context and the question) and each candidate separately, with a frozen Qwen3-8B and two projection heads, and scores them by cosine similarity, so the candidate embeddings can be cached.
Because the two are encoded apart, a candidate never attends to the context, and its training starts from tens of millions of pairs.
All of these start from a decoder.
We start from an encoder that reads the context, the question and all options together, and compare the two under the same training data.

\paragraph{Encoders that score options.}
Prompt-based classifiers use the masked-token head of an encoder with a verbalizer for a fixed set of labels \citep{schick2021pet, gao2021lmbff}.
Two encoder decision models are listed on public Jev leaderboards \citep{jevbench}: Laya \citep{laya} and Von \citep{von}, both built on ModernBERT-large with one mask per option.
Both add new layers on top of the encoder, and they are described in model cards.
Laya trains its head with reinforcement learning on proper scoring rules, and Von uses a listwise softmax loss with a Brier term.
We reuse the pretrained masked-token head and score each option by ``yes'' against ``no'', so there is no new head.
We also compare three training recipes, match the decoders in data and training, report a held-out test split, and add evidence.

\paragraph{Option letters, position and label bias.}
Reading the letter of the chosen option requires a model to bind the letter to the option text \citep{robinson2023mcsb}.
Models favour certain option labels, and their accuracy changes when the options are reordered \citep{zheng2024notrobust, pezeshkpour2024order}.
Corrections are usually applied after training, by calibrating with a content-free input \citep{zhao2021calibrate} or by dividing out a label prior and averaging over rotations of the options \citep{zhang2026anyjev}.
We shuffle the options during training and measure how the recipes differ in the labels they predict.

\paragraph{Encoders and decoders.}
Controlled studies pretrain encoders and decoders with the same recipe, and find that encoders do well on classification and retrieval \citep{weller2025ettin} and that masked language modelling suits tasks that need text representations \citep{gisserotboukhlef2025mlm}.
A much larger decoder can still beat an encoder on classification \citep{zhao2024classification}.
Decoder designs such as multi-query and grouped-query attention target long generation \citep{mqa, gqa}, and \citet{prefillonly} study what they cost when nothing is generated (\cref{sec:intro}).
Modern encoders keep long inputs and fast inference \citep{modernbert, ultrabert}.

\paragraph{Attribution and evidence.}
ContextCite learns a surrogate model from many ablations of the context \citep{contextcite}.
AT2 learns to weight attention heads, and needs a decoder pass \citep{at2}.
SelfCite fine-tunes a model to cite, using the change in the answer probability when the cited text is removed or kept alone as the reward \citep{selfcite}.
Rationale benchmarks such as ERASER compare selected evidence with human rationales \citep{deyoung2020eraser}.
We train an evidence head on the encoder with labels from human-marked evidence and from the SelfCite reward, which gives the evidence in one pass.
